\documentclass[11pt,a4paper]{article}

\usepackage[utf8]{inputenc}
\usepackage[T1]{fontenc}
\usepackage[brazil]{babel}
\usepackage{lmodern}
\usepackage{caption} 
\usepackage[left=2.5cm,right=2.5cm,top=2.5cm,bottom=2.5cm]{geometry}
\usepackage{graphicx}
\usepackage{xcolor}
\usepackage{listings}
\usepackage{amsmath,amssymb,amsthm}
\usepackage{float} 
 
\makeatletter
\newcommand{\inst}[1]{$^{#1}$}
\newcommand{\nextinstitute}{\par\vspace{2pt}}
\newcommand{\address}[1]{\gdef\@address{#1}}
\newcommand{\email}[1]{\gdef\@email{#1}}
\address{}
\email{}
\renewcommand{\maketitle}{%
  \thispagestyle{empty}
  \begin{center}
    {\Large\bfseries \@title \par}
    \vspace{0.9em}
    {\@author \par}
    \vspace{0.9em}
    {\small \@address \par}
    \vspace{0.4em}
    {\small\texttt{\@email} \par}
  \end{center}
  \vspace{1.2em}
}
\makeatother
% ==========================================
% CONFIGURAÇÃO DO CÓDIGO ISABELLE / ISAR
% ==========================================

\lstdefinelanguage{Isabelle}{
  keywords={
    theory, imports, begin, end, datatype, primrec, fun,
    definition, lemma, theorem, proof, qed, have, show,
    fix, assume, arbitrary, using, by, auto, simp, induction
  },
  keywordstyle=\color{blue}\bfseries,
  ndkeywords={cat, reverso, Nil, Cons},
  ndkeywordstyle=\color{purple}\bfseries,
  sensitive=true,
  comment=[l]{(*},
  morecomment=[s]{(*}{*)},
  commentstyle=\color{gray}\itshape,
  stringstyle=\color{red},
  basicstyle=\ttfamily\footnotesize,
  breaklines=true,
  breakatwhitespace=false,
  columns=fullflexible,
  keepspaces=true,
  numbers=left,
  numberstyle=\tiny\color{gray},
  frame=single,
  aboveskip=6pt,
  belowskip=6pt
}

\lstset{language=Isabelle}
 

\newcommand{\provaManuscrita}[3]{%
  \par\vspace{4pt}
  \centering
  \IfFileExists{#1}{%
    \includegraphics[
      width=0.92\textwidth,
      height=0.68\textheight,
      keepaspectratio
    ]{#1}%
  }{%
    \framebox{\parbox[c][5cm][c]{0.9\textwidth}{%
      \centering \color{gray}
      [ Imagem "\texttt{#1}" não encontrada ]
    }}%
  }
  \captionof{figure}{#2}
  \label{#3}
  \par\vspace{6pt}
}
\title{Métodos Formais --- Trabalho 1\\Provas por Indução de Especificações Equacionais Recursivas}
 
\author{Nicolas Borges Warmeling\inst{1},Germano Bruscato Corrêa\inst{1}}
 
\address{Pontifícia Universidade Católica do Rio Grande do Sul (PUCRS)\\
  90619-900 -- Porto Alegre -- RS -- Brasil\\
  Disciplina Métodos Formais -- Professor Júlio Machado\\}
 
\email{\{nicolas.warmeling,germano.bruscato\}@edu.pucrs.br}
 
\begin{document}
 
\maketitle
 
% ==========================================
% ESPECIFICAÇÃO DAS FUNÇÕES
% ==========================================
\section{Especificação das Funções}
 
Para a realização das demonstrações, foram consideradas as seguintes especificações de funções recursivas sobre o tipo indutivo de listas $List(\tau)$:
 
\begin{align*}
    &\text{cat}([], ys) = ys \tag{cateq1} \\
    &\text{cat}(x:xs, ys) = x:\text{cat}(xs, ys) \tag{cateq2} \\[6pt]
    &\text{reverso}([]) = [] \tag{reveq1} \\
    &\text{reverso}(x:xs) = \text{cat}(\text{reverso}(xs), [x]) \tag{reveq2}
\end{align*}
 
\vspace{0.3cm}
\hrule
\vspace{0.5cm}
 
% ==========================================
% LEMA 1: Associatividade da Concatenação
% ==========================================
\section{Lema Auxiliar 1: Associatividade de \texttt{cat}}
 
\noindent \textbf{Enunciado:}
\[ \forall xs, ys, zs \in List(\tau) \cdot \text{cat}(xs, \text{cat}(ys, zs)) = \text{cat}(\text{cat}(xs, ys), zs) \]
 
\subsection*{Prova Formal Feita à Mão}

\provaManuscrita{formais-1.png}{Demonstração detalhada e passo a passo à mão para o Lema 1.}{fig:lema1}
 
\subsection*{Prova em Isabelle / Isar}
\begin{lstlisting}[caption={Código-fonte Isar para o Lema 1}]
lemma l1: "\<forall>ys zs :: 'a list . cat xs (cat ys zs) = cat (cat xs ys) zs"
proof (induction xs)
  show "\<forall>ys zs :: 'a list . cat [] (cat ys zs) = cat (cat [] ys) zs"
  proof (intro allI)
    fix ys zs :: "'a list"
    have "cat [] (cat ys zs) = cat ys zs" by (simp only: cateq1)
    also have "... = cat (cat [] ys) zs" by (simp only: cateq1)
    finally show "cat [] (cat ys zs) = cat (cat [] ys) zs" by (simp)
  qed
next
  fix xs::"'a list" and x::'a
  assume HI:"\<forall> ys zs :: 'a list . cat xs (cat ys zs) = cat (cat xs ys) zs"
  show "\<forall>ys zs :: 'a list . cat (x#xs) (cat ys zs) = cat (cat (x#xs) ys) zs"
  proof (rule allI, rule allI)
    fix ys zs :: "'a list"
    have "cat (x#xs) (cat ys zs) = x # cat xs (cat ys zs)" by (simp only: cateq2)
    also have "... = x # cat (cat xs ys) zs" by (simp only: HI)
    also have "... = cat (x # cat xs ys) zs" by (simp only: cateq2)
    also have "... = cat (cat (x#xs) ys) zs" by (simp only: cateq2)
    finally show "cat (x#xs) (cat ys zs) = cat (cat (x#xs) ys) zs" by (simp)
  qed
qed
\end{lstlisting}

 
% ==========================================
% LEMA 2: Elemento Neutro à Direita
% ==========================================
\section{Lema Auxiliar 2: Elemento Neutro à Direita de \texttt{cat}}
 
\noindent \textbf{Enunciado:}
\[ \forall xs \in List(\tau) \cdot \text{cat}(xs, []) = xs \]
 
\subsection*{Prova Formal Feita à Mão}
\provaManuscrita{formais-2.png}{Demonstração detalhada e passo a passo à mão para o Lema 2.}{fig:lema2}
 
\subsection*{Prova em Isabelle / Isar}
\begin{lstlisting}[caption={Código-fonte Isar para o Lema 2}]
lemma l2: "cat xs [] = xs"
proof(induction xs)
  have "cat [] [] = []" by (simp only: cateq1)
  then show "cat [] [] = []" by (simp)
next
  fix xs::"'a list" and x::'a
  assume HI: "cat xs [] = xs"
  have "cat (x#xs) [] = x # cat xs []" by (simp only: cateq2)
  have "x # cat xs [] = x # xs" by (simp only: HI)
  then show "cat (x#xs) [] = x # xs" by (simp)
qed
\end{lstlisting}
 

 
% ==========================================
% LEMA 3: Reverso da Concatenação
% ==========================================
\section{Lema Auxiliar 3: Distribuição do \texttt{reverso} sobre \texttt{cat}}
 
\noindent \textbf{Enunciado:}
\[ \forall xs, ys \in List(\tau) \cdot \text{reverso}(\text{cat}(xs, ys)) = \text{cat}(\text{reverso}(ys), \text{reverso}(xs)) \]
 
\subsection*{Prova Formal Feita à Mão}
\provaManuscrita{formais-31.png}{}{fig:lema3.1}
\provaManuscrita{formais-32.png}{Demonstração detalhada e passo a passo à mão para o Lema 3.}{fig:lema3.2}
 
\subsection*{Prova em Isabelle / Isar}
\begin{lstlisting}[caption={Código-fonte Isar para o Lema 3}]
lemma l3: "\<forall>ys :: 'a list . reverso (cat xs ys) = cat (reverso ys) (reverso xs)"
  proof (induction xs)
    show "\<forall>ys :: 'a list . reverso (cat [] ys) = cat (reverso ys) (reverso [])"
    proof (intro allI)
      fix ys :: "'a list"
      have "reverso (cat [] ys) = reverso ys " by (simp only: cateq1)
      also have "... = cat (reverso ys) [] " by (simp only: l2)
      also have "... = cat (reverso ys) (reverso [])" by (simp only: reveq1)
      finally show "reverso (cat [] ys) = cat (reverso ys) (reverso [])" by (simp)
    qed
  next
    fix xs :: "'a list" and x::'a
    assume HI: "\<forall> ys :: 'a list . reverso (cat xs ys) = cat (reverso ys) (reverso xs)"
    show "\<forall>ys :: 'a list . reverso (cat (x#xs) ys) = cat (reverso ys) (reverso (x#xs))"
    proof (intro allI)
      fix ys :: "'a list"
      have "reverso (cat (x#xs) ys) = reverso (x#(cat xs ys)) " by (simp only: cateq2)
      also have "... = cat (reverso (cat xs ys)) [x]" by (simp only: reveq2)
      also have "... = cat (cat (reverso ys) (reverso xs)) [x]" by (simp only: HI)
      also have "... = cat (reverso ys) (cat (reverso xs) [x])" by (simp only: l1)
      also have "... = cat (reverso ys) (reverso (x#xs))" by (simp only: reveq2)
      finally show "reverso (cat (x#xs) ys) = cat (reverso ys) (reverso (x#xs))" by (simp)
    qed
  qed
\end{lstlisting}
 
\newpage
 
% ==========================================
% TEOREMA PRINCIPAL: Involução do Reverso
% ==========================================
\section{Teorema Principal: Involução do \texttt{reverso}}
 
\noindent \textbf{Enunciado:}
\[ \forall xs \in List(\tau) \cdot \text{reverso}(\text{reverso}(xs)) = xs \]
 
\subsection*{Prova Formal Feita à Mão}
\provaManuscrita{formais-41.png}{}{fig:teoremapt1}
\provaManuscrita{formais-42.png}{Demonstração detalhada e passo a passo à mão para o Teorema Principal.}{fig:teoremapt2}
 
\subsection*{Prova em Isabelle / Isar}
\begin{lstlisting}[caption={Código-fonte Isar para o Teorema Principal}]
theorem t1: "reverso (reverso xs) = xs"
  proof (induction xs)
    have "reverso (reverso []) = reverso []" by (simp only: reveq1)
    also have "reverso [] = []" by (simp only: reveq1)
    then show "reverso (reverso []) = []" by (simp)
  next
    fix xs :: "'a list" and x::'a
    assume HI: "reverso (reverso xs) = xs"
    have "reverso (reverso (x#xs)) = reverso (cat (reverso xs) [x])" by (simp only: reveq2)
    also have "... = cat (reverso [x]) (reverso (reverso xs))" by (simp only: l3)
    also have "... = cat (reverso (x#[])) xs" by (simp only: HI) (* transf. [x] \<rightarrow> x#[] nesse passo *)
    also have "... = cat (cat (reverso []) [x]) xs" by (simp only: reveq2)
    also have "... = cat (cat [] [x]) xs" by (simp only: reveq1)
    also have "... = cat (x#[]) xs" by (simp only: cateq1) (* transf. [x] \<rightarrow> x#[] nesse passo *)
    also have "... = x # (cat [] xs)" by (simp only: cateq2)
    also have "... = x # xs" by (simp only: cateq1)
    finally show "reverso (reverso (x#xs)) = x # xs" by (simp)
  qed
\end{lstlisting}
 
\end{document}
 

